% Erklaerung der Autoren: Version 1.0 (2021, als Verweis) und Version 2.0.0 (neu).
% ANNAHME: Die Erklaerung zu 2.0.0 ist fuer Tim M. Kretzschmar formuliert (Autor der
% Commits auf dem Branch v2.0.0). Falls Konstantin Blechschmidt ebenfalls mitunterzeichnen
% soll, unten den zweiten Unterschriftsblock einkommentieren. Die Formulierung zur
% KI-Unterstuetzung kann jederzeit angepasst werden.
\section{Erklärung der Autoren}\label{sec:erklaerung}

\subsection{Zur Version 1.0 (Abgabe 2021) -- Verweis}
Zur Abgabe im Jahr 2021 (Version~1.0, Stand im Git: \commit{89f8212}) haben beide Autoren am 03.08.2021 eine eidesstattliche Erklärung abgegeben. Sie bezieht sich ausschließlich auf die damalige Fassung des Programms und der Dokumentation. Diese Erklärung ist hier unverändert wiedergegeben, damit nachvollziehbar bleibt, welche Fassung sie betrifft; sie ersetzt \emph{nicht} die Erklärung zu Version~\versionProjekt\ im folgenden Abschnitt.

\begin{quote}
\emph{Hiermit bestätigen wir, Konstantin Blechschmidt \& Tim M.\ Kretzschmar, dass wir dieses Programm \glqq AES -- in TASM umgesetzt\grqq{} selbstständig und nur unter Angabe der genutzten Quellen erstellt und getestet haben.}
\end{quote}

\vspace{1mm}
\noindent
\begin{minipage}{\textwidth}
    \centering
    \small\textbf{Unterschriften zur Version 1.0 (Stand 03.08.2021)}\par\vspace{2mm}
    \begin{minipage}[b]{0.52\textwidth}
        \includegraphics[width=0.8\textwidth]{images/unterschrift_tim.png}
    \end{minipage}\hfill
    \begin{minipage}[b]{0.44\textwidth}
        \raggedright Tim M.\ Kretzschmar\\(Wermsdorf, 03.08.2021)
    \end{minipage}

    \vspace{4mm}

    \begin{minipage}[b]{0.52\textwidth}
        \includegraphics[width=0.7\textwidth]{images/unterschrift_konstantin.png}
    \end{minipage}\hfill
    \begin{minipage}[b]{0.44\textwidth}
        \raggedright Konstantin Blechschmidt\\(Schkeuditz, 03.08.2021)
    \end{minipage}
\end{minipage}\par

%------------------------------------------------------------------------------
\Needspace{22\baselineskip}
\subsection{Zur Version \versionProjekt}\label{sec:erklaerung-v2}
Die Version~\versionProjekt\ ist eine Weiterentwicklung der Version~1.0. Die Änderungen sind in Kapitel~\ref{sec:aenderungen} beschrieben und im Git-Branch \texttt{v2.0.0} einzeln nachvollziehbar.

\begin{quote}
Hiermit erkläre ich, Tim M.\ Kretzschmar, dass ich die Weiterentwicklung \glqq AES -- in TASM umgesetzt\grqq{} zur Version~\versionProjekt\ verantworte und die genutzten Quellen angegeben habe. Bei der Fehlersuche, der Korrektur des Quelltextes, der Erstellung des Python-Modelltests (\asm{Dokumentation/Modelltest/modelltest.py}) und der Überarbeitung dieser Dokumentation wurde die KI \emph{Claude} (Anthropic) unterstützend eingesetzt; sie ist in den entsprechenden Git-Commits als Co-Autor eingetragen. Die Ergebnisse habe ich geprüft und übernommen.
\end{quote}

\textbf{Stand der Prüfung.} Die Korrekturen der Version~\versionProjekt\ wurden an einem Python-Nachbau der Assemblerroutinen geprüft (Abschnitt~\ref{sec:tests-modell}). Der Lauf des übersetzten Programms in der DOS-Box steht noch aus (Abschnitt~\ref{sec:tests-dosbox}); die Erklärung in der vorliegenden Form bezieht sich daher auf die Erstellung, nicht auf einen abgeschlossenen DOS-Box-Test.

\vspace{8mm}
\noindent
\begin{minipage}{\textwidth}
    \begin{minipage}[b]{0.46\textwidth}
        \rule{\textwidth}{0.4pt}\\[1mm]
        \small Ort, Datum
    \end{minipage}\hfill
    \begin{minipage}[b]{0.46\textwidth}
        \rule{\textwidth}{0.4pt}\\[1mm]
        \small Tim M.\ Kretzschmar
    \end{minipage}
\end{minipage}\par

% Zweiter Unterschriftsblock fuer Konstantin Blechschmidt (nur falls er die
% Weiterentwicklung ebenfalls verantwortet):
% \vspace{10mm}
% \noindent
% \begin{minipage}{\textwidth}
%     \begin{minipage}[b]{0.46\textwidth}
%         \rule{\textwidth}{0.4pt}\\[1mm]
%         \small Ort, Datum
%     \end{minipage}\hfill
%     \begin{minipage}[b]{0.46\textwidth}
%         \rule{\textwidth}{0.4pt}\\[1mm]
%         \small Konstantin Blechschmidt
%     \end{minipage}
% \end{minipage}\par
